\section{Material Girls and Spiritual Tourists}

\begin{quotex}
Many individuals, having reached a certain contemplation of the mystery, feel rather like shooting themselves in the head or returning to the pleasures of life, because this game and this trip is something that does not lead anywhere, or which has no meaning.
\flright{\textsc{Ghislaine Gualdi}}
\end{quotex}


These personality studies are not intended to denigrate anyone but rather to illustrate the human condition and perhaps illuminate a moral point. Some did not make the cut. There was one woman who used to attend my seminars on non-duality; she claimed she could see the ghosts of departed Indian braves in the empty seats. I never found out if they were just following her around or were genuinely interested in non-duality.

Then there were the two women who kept a dossier on me for several months, including my emails, lecture notes, and even data on my friends and acquaintances. It had reached more than 100 pages by the time I found out about it.

Just this week I found some 20 year old transcriptions of material channeled by a Swiss woman, Ghislaine Gualdi, from somewhere high up on the Chain of Being. I had even witnessed in person some of her private talks at which she would be posed a question to which she responded from a trance state. I don't know what happened to her, but her material is too nihilistic for popular consumption. At the end of the transcribed session, the spirit being tells us:

\begin{quotex}
I do not tell you to change your life starting tomorrow because of what I have told you. You will not be able to.
\end{quotex}


The characters described below will not be able to change their lives, and neither will you. However, you may be able to decide how to deal with that realization.

\paragraph{The Tragedy of Life}


\begin{quotex}
The fact that there are many ignorant people speaking on the earth is a tragedy.
\flright{\textsc{Ghislaine Gualdi}}
\end{quotex}


\begin{quotex}
Each Personality has a principal feature for its axis, and all its qualities and defects gravitate around this. It is not necessary for this trait to be striking; it can be insignificant and even ridiculous.
\flright{\textsc{Boris Mouravieff}}
\end{quotex}


\textbf{Aristotle} recognized this chief feature in Greek tragedies as the flaw in the protagonist that leads to his downfall. He used the word \emph{hamartia\footnote{\url{https://en.wikipedia.org/wiki/Hamartia}}} to describe it, which literally means to ``miss the mark''. Not by chance is the same word used to denote ``sin'', which is the tragedy of human existence that misses the mark of \emph{theosis}. Specifically, however, it is not sin in general, but rather the unique fault, or even illusion, that is the axis around which all other faults move.

Curiously, people can notice the chief features of everyone else except their own. That is the key to the drama: the audience can see it all unfold yet the protagonist is taken by surprise despite his central role in bringing about his own ruin. They just can't prevent it from happening.

\paragraph{The Material Girl}


If the home is a symbol of the womb and the automobile a phallic symbol, then Magda was condemned to relive the Freudian drama. A successful medical researcher, Magda visited her flat in a city in Budapest each summer. Hopeful to discuss old Europe, or even Jobbik, I picked her up at her townhouse. It was in immaculate condition as though no one was even living there. In effect, it was like a house with a still intact hymen.

Since it was a nice night, we took a ride to the beach in my Sonata to get a gelato at the gelateria owned by an older Sicilian couple. She was more interested in talking about the size and the assessed value of her previous house in New Jersey, and her plans to upgrade in Florida. She wanted to buy the larger unit next door, but complained about its owner who was a ``hoarder'', with ``stuff'' all over the place. I immediately thought of my own environs with a couple of thousand books strewn around in shelves, boxes, on the floor, tables, and even sofas: I made a mental note to keep her away.

Eventually, I did manage to bring the conversation around to more ``spiritual'' topics. She did feign some curiosity about Gornahoor and the Gnosis study group, although it is really not something I can describe easily on an elevator ride. Nevertheless, since there is an aim for women to participate, I persisted.

When I described the typical demographic as younger men in revolt against the modern world, she suddenly became animated. ``I see,'' she said, ``It is for men who are threatened by powerful women like me.'' That did take me by surprise despite efforts at external considering. However, it is not far off. It is not her as an individual that is fearsome, but rather the archetype of the Demetrian force that she represented. That may be the chief feature of the modern world.

The next morning she sent me this text: ``I prefer expensive cars''.

\paragraph{The Alien Tourist}


Catching up on movies that I've missed, I watched about 20 minutes of \emph{Aliens}. The female protagonist, Ripley, was both gentle with a young girl and decisive around the men, not to mention intelligent. In the scene I watched, the captain became paralyzed by fear and indecisiveness. Ripley overpowered him and took control of the vehicle to save the remnant. She then dominated not only the next in command but also the businessman who was only interested in preserving his property. No one has that many skills (and others I've skipped over).

I mentioned this to a successful woman, a business consultant with a Harvard MBA. She loved that movie and I guess I know why. When I mentioned my objections, she just advised me to ``suspend reality'' and enjoy it. Now it is one thing, I suppose, to suspend reality for a movie, but we are being asked to extend that suspension to life itself.

\paragraph{The Artificial Tourist}


I also saw part of \emph{Artificial Intelligence} by \textbf{Steven Spielberg}. Now Mr Spielberg is obviously toying with us, demonstrating how an artificial world can generate real events. His first attempt, years ago, was with \emph{ET}, which enticed the audience to empathize with a doll. Traditionally, drama was intended to expose the human condition and derive moral lessons from it. With ET, there is no human condition and no morality; rather it is an experiment to manipulate human emotion. If the audience feels sorrow for the misfortunes of the doll or joy for its successes, moral concerns must recede. There is no ``tragic flaw'' or hamartia in the artificial protagonist; rather, it is always the fault of someone else.

AI brings ET to its logical conclusion through the creation of a golem\footnote{\url{https://en.wikipedia.org/wiki/Golem}}. An android which simulates human emotion is equivalent to, if not actually superior to, an actual human. I've heard that many cried at the movie's ending, which makes as much sense as bawling when your electric can opener craps out.

The point is that people today lead a factitious existence, barely able to distinguish reality from their illusions.

\paragraph{The Gleeful Tourist}


Elaine was jabbering incessantly at a suburban barbeque, at which most of the guests were deracinated, ``spiritual but not religious'', types. Elaine , on the other hand, was tribal, theistic, and religiously committed to her very ancient monotheistic faith. While discussing her recent European trip, she mentioned her visit to Prague. This included trips to its many churches, which are little more than museums now. With a great amount of glee, she informed us of the fact, explaining --- correctly --- that 80\% of Czechs are religiously uncommitted.

In her joy at the elimination of Christianity in the center of Europe, she makes common cause with various neopagan, new right, secular right types, even if ostensibly for opposite reasons.

\paragraph{The Spiritual Tourist}


I met Mary at the café in a bookstore. Although an Irish Catholic, with a parochial school education, her commitment to her birth religion was weak. She got a law degree, but gave up active practice to raise a family. Her ex-husband was a high powered attorney who recently ran off with a much younger Puerto Rican woman. Although she was eager to tell it, I wasn't really interested in that story.

Now 60 years old, her family gone, no career, but plenty of money, Mary was looking for something to kill time. When confronted with the stark realization that her life trip did not really lead anywhere, she endeavored to sink into the pleasures of life. She took up hobbies that she really should have done when she was young. She joined a private dance studio, spending thousands of dollars per year on lessons; it hadn't quite dawned on her that her only dance partner in life would always be her instructor. She also started yoga classes. Several sea cruises per year killed time for her. Her three pound micro-dog, named Jacqueline Kennedy Onassis, seemed to provide companionship when she was at home.

I was able to explain Tradition in a way that interested her as well as methods to master one's own mind. Knowing nothing of Buddhism, she and her yoga instructor had recently gone on a retreat to a Buddhist monastery in Nepal. She mentioned her attempts at meditation and I offered to teach her if she wanted.

She told me she wanted to read \emph{The Power of Now} by Eckhart Tolle, since it had been recommended by a friend. Suddenly, something awoke in her and she exclaimed, ``But I could just ask you about stopping thoughts!'' Since our coffees had long since gotten cold, we decided to part ways. I saw her wandering toward the bookshelves in search of her book, desperately hoping for her next pointless distraction.


\flrightit{Posted on 2015-08-25 by Cologero}

\begin{center}* * *\end{center}

\begin{footnotesize}
\begin{sffamily}
\texttt{White Rabbit on 2015-08-26 at 16:35 said:}

``Curiously, people can notice the chief features of everyone else except their own. That is the key to the drama: the audience can see it all unfold yet the protagonist is taken by surprise despite his central role in bringing about his own ruin.''

This reminded me of a recommendation I read awhile ago and forgot from Franz Bardon:

``Arrange for a magic diary and enter all the bad sides of your soul into it. This diary is for your own use only, and must not be shown to anybody else. It represents the so-called control book for you. In the self-control of your failures, habits, passions, instincts and other ugly character traits, you have to observe a hard and severe attitude towards yourself.

Be merciless towards yourself and do not embellish any of your failures and deficiencies. Think about yourself in quiet meditation, put yourself back into different situations of your past and remember how you behaved then and what mistakes or failures occurred in the various situations. Make notes of all your weaknesses, down to the finest nuances and variations. The more you are discovering, all the better for you.

Some especially endowed disciples have been able to discover hundreds of failures in the finest shades... This psychic preliminary work is indispensable to obtain the magic equilibrium, and without it, there is no regular progress of the development to be thought of.''

\url{http://metapara.weebly.com/1-introspection-or-self-knowledge-ndashhermeticism.html}


That way, you can read your own drama. But you might need to buy yourself a safe!

\hfill

\texttt{White Rabbit on 2015-09-10 at 02:27 said:}

About AI and illusions... now, I am an absolute brutal fascist when it comes to our automated friends. I think that you can beat, enslave, and torture your little humanoid can openers to your heart's content (leaving the question of your soul's health out of it). Alas, if you were to ask the men (and mechanisms) of the future, they'd tell you that this makes me a horrible person, perhaps even subhuman. The idea of ``robot rights'' has already crept into the mainstream.

\url{https://en.wikipedia.org/wiki/Ethics\_of\_artificial\_intelligence\#Robot\_rights}


Did I mention that I need to buy a safe and a diary? I can't forget, or else I'll be forever trapped down in the mud with these phantom-chasing lunatics!

\hfill

\texttt{Matt on 2015-09-10 at 20:16 said:}

White Rabbit,

Not just robots. There are also those that believe rivers and trees have inalienable rights and that these rights should be recognized by law. It seems at no point have they considered just how exactly these things are able to assert their apparent rights; nor have they considered how these things could be held to account for failing to live up to their duties and obligations (though I guess it doesn't matter for the liberal mind, as rights are divorced from obligations).


\hfill

\end{sffamily}
\end{footnotesize}

